\documentclass[11pt]{article}

\usepackage{amsmath}
\usepackage{amsfonts}
\usepackage{amsthm,thmtools}
\usepackage{tikz}
\usepackage{enumerate}


\newcommand{\batsym}{\begin{tikzpicture}[scale=0.05]
	\draw[fill,black]   (-0.25,1.48) .. controls (-0.1,1.5) and (0.1,1.5) .. (0.25,1.48) -- (0.35,1.92) .. controls (0.425,1.8) and (0.41,1.3) .. (0.45,1.2) .. controls (0.6,1.05) and (1.96,1.05) .. (1.98,2.08) -- (5.93,2.08) .. controls (4.2,1.45) and (4,0.3) .. (4.2,-0.28) .. controls (2.4,-0.09) and (0.4,-0.5) .. (0,-2.052) .. controls (-0.4,-0.5) and (-2.4,-0.09) .. (-4.2,-0.28) .. controls (-4,0.3) and (-4.2,1.45) .. (-5.93,2.08) -- (-1.98,2.08) .. controls (-1.96,1.05) and (-0.6,1.05) .. (-0.45,1.2) .. controls (-0.41,1.3) and (-0.425,1.8) .. (-0.35,1.92) -- (-0.25,1.48);
	\end{tikzpicture}}
\newcommand{\QEB}{\begin{flushright}\batsym\end{flushright}}



\def\bold{\textbf}
\theoremstyle{plain}
\newtheorem{problem}{Problem}
\theoremstyle{remark}
\newtheorem*{idea}{Idea}

\renewcommand*{\proofname}{Solution}


\renewcommand{\qedsymbol}{$\batsym$}

\newcommand {\ds}{\displaystyle\sum}
\newcommand {\nnR} {\mathbb{R}_{\geq 0}}

\newcommand {\Mtt}[4] {\begin{bmatrix}#1 & #2\\ #3 & #4\end{bmatrix}}
\newcommand {\Ctt}[2] {\Mtt {#1}{-#2}{#2}{#1}}
\newcommand {\mbR} {\mathbb R}
\newcommand {\mcC} {\mathcal C}
\newcommand {\mbZ} {\mathbb Z}
\newcommand {\mbN} {\mathbb N}
\newcommand {\mbQ} {\mathbb Q}
\newcommand {\RnQ} {\mathbb R \setminus \mathbb Q}
\newcommand {\mcI} {\mathcal I}
\newcommand {\dlim}[2] {\displaystyle \lim_{#1 \rightarrow #2}}

\newcommand {\eps} {\varepsilon}
\newcommand {\del} {\delta}
\newcommand {\abs}[1]{\left| #1 \right|}

\newcommand {\pwii}[6] {\begin{cases} 
		#1 & #2 \\
		#3 & #4 \\
		#5 & #6
\end{cases}}

\begin{document}
	\title{Math 3A03 - Tutorial 8 Solutions - Winter 2019}
	\author{Nikolay Hristov}
	\date{March 11/13, 2019}
	\maketitle
	
	\begin{problem}
		Let \[f(x) = \pwii{\frac 1 q}{x=\frac p q,p,q\in\mbN, \quad\gcd(p,q)=1}{0}{x\notin \mbQ}{0}{x=0}.\] For which $y\in \mbR$ does the limit exist? Where is this function continuous?
	\end{problem}

\begin{proof}
	The limit exists and the function is continuous when $y\in \mbR\setminus \mbQ$, and it does not otherwise. Notice: this is a solution to one of your challenge problems in class! The function is often called Thomae's function if you'd like to read more about it.
	
	For now let's prove that it is continuous at each irrational number and not continuous at each rational number, or at least sketch a proof.
	
	First we'll show that if $y\in \mbQ$ then \[\dlim x y f(x)\] does not exist. Since y is rational we have that $f(y)=\frac 1 q$ for some $q\in \mbN$. Fix $\eps = \frac 1 q$, then given any $\del>0$ notice that by the density of the irrational numbers there exists an irrational $x\in (y-\del,y+\del)$, hence $f(x)=0$, so we have an $x$ for which $\abs{x-y}<\del$ but $\abs{f(x)-f(y)}=\abs{f(y)}=\frac 1 q = \eps\geq \eps$, hence the limit does not exist.
	
	Now suppose that $y\notin \mbQ$, we'd like to show that \[\dlim x y f(x)=f(y)=0.\] Further suppose that $y\in [0,1]$, this wasn't given as part of the function, but we can do this by without loss of generality since for any $x\in \mbR$, we have $y\in [0,1], n\in \mbN$ with the property that $x=y+n$ and $f(y+n)=f(y)$. Given $\eps>0$ we know (by the Archimedean property) that there exists an $n\in \mbN$ such that $\frac 1 {n+1} < \eps$. Our goal here is to find an interval $(y-\del,y+\del)$ such that $\abs{f(x)-f(y)}=\abs {f(y)}\leq \frac 1 {n+1} < \eps$. Which is to say that each rational number $x$ in that interval when written as $x=\frac p q, \gcd(p,q)=1$ has the property that $n+1\leq q.$ Define the set $S=\{\frac p q: 1\leq p\leq q \leq n,\gcd(p,q)=1\}$, this set contains all rational numbers between 0 and 1 with denominator $q\leq n.$
	The cardinality of this set is at most $|S|\leq \ds_{k=1}^n k = \frac {n(n+1)} 2,$ this is a finite set. Define $d=\inf\{|y-s|:s\in S\}$, since $S$ is a finite set and $y$ is irrational this value exists and is attained for some $s^\star \in S$ and $d>0$. Let $\del = \frac d 2$, then for each $x\in [y-\del,y+\del]$ we don't have any rational numbers with denominators $q\leq n$, hence if $|x-y|<\del$ then $\abs{f(x)-f(y)}=\abs{f(x)}\leq \frac 1 {n+1}<\eps$.
\end{proof}

	\begin{problem}
		Which of the following are uniformly continuous:
		
		\begin{enumerate}
			\item $f(x)=x^3$ with $x\in\mbR$
			\item $f(x)=x^3$ with $x\in [0,3]$
			\item $f(x)=\frac{1}{x}$ with $x\in(0,1)$
			\item $f(x)=\frac 1 x$ with $x\in [1,\infty)$
			\item $f(x)=\sin\left(\frac 1 x\right)$ with $x\in (0,1)$
		\end{enumerate}
	\end{problem}

	\begin{proof}
	2 and 4 are uniformly continuous. 
	\begin{enumerate}
		\item We'll prove this by contradiction. Suppose that given $\eps=1>0$, there is a $\del>0$ such that for any $x,y\in \mbR$ such that if $|x-y|<\del$ (in particular if $|x-y|=\frac \del 2<\del$) then $\eps>|x^3-y^3|=|x-y||x^2+xy+y^2|=\frac \del 2 |x^2+xy+y^2|$. But if $x>y\geq \frac 1 {\sqrt{\del}}$ we have $1>|x^3-y^3|=\frac \del 2 |x^2+xy+y^2| > \frac \del 2 \frac 3 \del = \frac 3 2,$ which is a contradiction.
		\item $x^3$ is continuous, and $[0,3]$ is a compact set, hence $x^3$ is uniformly continuous on the set.		
		\item We'll prove this by contradiction, suppose that $\frac 1 x$ is uniformly continuous, then there is a $\delta>0$ such that for any $x,y\in(0,1)$, if $|x-y|<\delta$ then $|\frac 1 x - \frac 1 y|<1$. Note in particular this means that if $|x-y|=\frac \delta 2 < \delta$ then $\frac{\delta}{2xy}=|\frac{1}{x}-\frac{1}{y}|<1$. Since $0<x<1$ then $\frac 1 y<\frac 1 {xy},$ hence $\frac{\delta}{2y}<1.$ If $\delta$ is sufficiently small then $0<\frac \delta 2<1$, but since a smaller $\del$ always works we can wlog assume $\del <2$ so we can let $y=\frac \delta 2$, in which case we get that $1<|\frac{1}{x}-\frac{1}{y}|<1$, which is a contradiction. 
		\item Given $\epsilon>0$, let $\delta=\epsilon$. since $x,y\geq 1$, meaning $0<\frac{1}{x}\leq 1,0<\frac 1 y \leq 1$. Suppose that $|x-y|<\delta$, we have $|\frac{1}{x}-\frac 1 y|=\frac{|y-x|}{|xy|}\leq |y-x|<\delta=\epsilon$. Since this holds for any $x,y$ on $[1,\infty)$ the function is uniformly continuous.
		\item Hint: consider the sequence of points $x_n=\frac{2}{\pi(2n+1)},y_n=\frac{1}{2\pi n}$, show that given $\delta>0$ there is an $N\in \mbN$ so that $|x_n-y_n|<\delta$ for all $n\geq N$, and that $|f(x_n)-f(y_n)|=1$. What does this tell you about the uniform continuity of $\sin(1/x)$ on the specified interval?
	\end{enumerate}
\end{proof}



\begin{problem}
	Let $f:\mbR \to \mbR$ be a continuous function with the property that 
	\begin{equation}
	\label{limits}
	\lim_{x\to -\infty} f(x) = \lim_{x\to \infty} f(x) = 0.
	\end{equation}
	Show that $f$ has either an absolute maximum or an absolute minimum but not necessary both.
\end{problem}



\begin{proof}
	
	Next week!
	
	\iffalse
	We assume that the function $f$ is not identically zero on the real line, otherwise both absolute maximum and minimum of $f$ is zero.
	
	Since $f$ is not identically zero on the real line, there exists a point $p \in \mathbb{R}$ satisfying $f(p) \neq 0$. Set $\epsilon:= |f(p)|>0$. 
	
	Consider the set $S:=\{x: |f(x)| \geq \epsilon \}$. Note that the set $S$ is bounded from below. Otherwise, there exists a sequence $\{x_n\} \subseteq S$ such that $x_n \to -\infty$ and $f(x_n) \to \epsilon$, which contradict to the condition (\ref{limits}) since $\epsilon>0$. Similarly, we prove that $S$ is also bounded from above. Therefore, both infimum and supremum of $S$ exist. 
	
	Consider any accumulation point of the set $S$, i.e. consider a sequence 
	$\{x_n\} \subseteq S$ such that $x_n \to x \in \mbR$. Since $f$ is continuous, $f(x_n) \to f(x)$. Note that for each $n \in \mbN$, $f(x_n) \geq \epsilon$, and so $f(x) \geq \epsilon$ also. Then, $x \in S$, which means that the set $S$ is closed. Therefore, inf$(S) \in S$ and sup$(S) \in S$.
	Set $c_1:= {\rm inf}(S)$ and $c_2:= {\rm sup}(S)$.
	
	By the above construction, 
	$$
	|f(x)|<\epsilon = |f(c_1)| \quad {\rm for} \ {\rm all} \quad x<c_1, \quad {\rm and}
	$$
	$$
	|f(x)|<\epsilon = |f(c_2)| \quad {\rm for} \ {\rm all} \quad x>c_2.
	$$
	Therefore, $f(c_1)$ is either an absolute maximum or an absolute minimum of $f$ on the interval $(-\infty, c_1]$, and $f(c_2)$ is either an absolute maximum or an absolute minimum of $f$ on the interval 
	$[c_2, +\infty)$.
	
	If $c_1 = c_2$, then $f(c_1)$ is an absolute extremum of $f$ on $\mbR$. If $c_1 \neq c_2$, then the interval $[c_1,c_2]$ is nonempty, and by Extreme Value Theorem, there exist the absolute maximum $M$ and the absolute minimum $m$ of $f$ on the interval $[c_1, c_2]$. 
	If max$\big(f(c_1),f(c_2)\big)>0$, then max$\big(f(c_1),f(c_2), M\big)$ is an absolute maximum of $f$ on $\mbR$. If min$\big(f(c_1),f(c_2)\big)<0$, then min$\big(f(c_1),f(c_2), m\big)$ is an absolute minimum of $f$ on $\mbR$. If 
	$$
	{\rm max} \big(f(c_1),f(c_2)\big)= {\rm min} \big(f(c_1),f(c_2)\big)=0,
	$$ then max$\big(f(c_1),f(c_2), M\big)$ is an absolute max of $f$ on $\mbR$ and min$\big(f(c_1),f(c_2), m\big)$ is an absolute min of $f$ on $\mbR$.
	
	There exist cases when there is no absolute max or absolute min. For example, the function $f(x) = \exp(-x^2)$ does not have absolute min, whereas the function $f(x) = -\exp(-x^2)$ does not have absolute max.
	
	\fi
\end{proof}


	\begin{problem}
	Let $f$ be a continuous, one-to-one function defined on the interval $[a,b]$ with $f(a)<f(b)$. Show that, for all $x,y \in [a,b]$, if $x<y$ then $f(x)<f(y)$.
\end{problem}
	
	
	\begin{proof}
		Next week!
		\iffalse
		Assume, by contradiction, that there exist $x,y \in [a,b]$ such that $x<y$ and $f(x)>f(y)$. Note that $a \leq x <y \leq b$, and so $a\neq y$. 
		We divide the proof into three cases. 
		
		{\bf Case 1:} $f(y) = f(a)$. Then, since $a\neq y$, we get a contradiction to the injectivity of $f$.
		
		{\bf Case 2:} $f(y) < f(a)$. We consider two intervals $(a,y)$ and $(y,b)$. The second interval is nonempty since $f(y)<f(a)<f(b)$, and so $y < b$. Then, for the interval $(a,y)$, choose any $d \in (f(y), f(a))$. By Intermediate Value Theorem, there exist some $c_1 \in (a,y)$ such that $f(c_1) = d$. For the interval $(y,b)$, since $f(y)<f(a)<f(b)$ holds, 
		our previously chosen $d \in (f(y), f(b))$. Then, by Intermediate Value Theorem, there exist some $c_2 \in (y,b)$ such that $f(c_2) = d$. Since $(a,y) \cap (y,b) = \emptyset$, $c_1 \neq c_2$. Therefore, $f(c_1) = d = f(c_2)$ contradicts to the injectivity of $f$.
		
		{\bf Case 3:} $f(y) > f(a)$. We consider two intervals $(a,x)$ and $(x,y)$. The first interval is nonempty since $f(a)<f(y)<f(x)$, and so $a < x$. Then, for the interval $(a,x)$, choose any $d \in (f(y), f(x)) \subset (f(a), f(x))$. By Intermediate Value Theorem, there exist some $c_1 \in (a,x)$ such that $f(c_1) = d$. For the interval $(x,y)$, our previously chosen 
		$d \in (f(y), f(x))$. Then, by Intermediate Value Theorem, there exist some $c_2 \in (x,y)$ such that $f(c_2) = d$. Since $(a,x) \cap (x,y) = \emptyset$, $c_1 \neq c_2$. Therefore, $f(c_1) = d = f(c_2)$ contradicts to the injectivity of $f$.
		
		In each case we have a contradiction. Therefore, for all $x,y \in [a,b]$, if $x<y$ then $f(x)<f(y)$.
		
		\fi
	\end{proof}
\end{document}